\renewcommand{\thetable}{3.A.13}
\begin{tablalafrox}{Criterios de aceptación. Fuente: elaboración propia a partir de los resultados de aceptación exigidos por Distribuidora Puelche, con metas propuestas por LafroX donde se indica.}%
  {tab:3-A-13}%
  {F{0.065} F{0.216} F{0.126} F{0.168} Y}%
  {\cab{ID} & \cab{Resultado exigido} & \cab{Situación actual} & \cab{Meta} & \cab{Momento / método}}
  R18-01 & Lista de clientes afectados por lote con evidencia en menos de dos horas. & 9 días, con resultado estimado. & < 2 horas. & Marcha blanca E1: simulacro de retiro con un lote real, cronometrado. \\
  R18-02 & El 100 \% de las recepciones registra el lote de los productos que lo requieren. & 41 \% sin registro en el producto involucrado. & 100 \%. & Marcha blanca E1: recepciones con lote sobre recepciones que lo requieren, cuatro semanas. \\
  R18-03 & Registro continuo y automático de temperatura en cámaras y vehículos, auditable y disponible para el cliente. & 3 lecturas manuales por viaje. & 100 \% de cámaras y vehículos con frío con registro continuo y automático, auditable y disponible para el cliente. & Marcha blanca E1: auditoría de series térmicas sin brechas y comprobación de que el CLIENTE puede consultar los registros e identificar la cámara o vehículo y la fecha y hora de cada lectura. \\
  R18-04 & OTIF medido de una sola forma y en la meta comprometida. & 82,4 \%, medido de forma distinta por cada área. & 90 \% en el mes 15; 93 \% en el mes 19; 95 \% en el mes 32 (tramos propuestos por LafroX; 95 \% corresponde a la meta de referencia de Distribuidora Puelche). & Medición diaria con RNG-09; revisión mensual en comité. \\
  R18-05 & El preventista conoce stock y crédito al tomar el pedido. & No los conoce. & 100 \% de pedidos con stock y crédito visibles. & Marcha blanca E1: registro de consultas por pedido. \\
  R18-06 & Ningún pedido se pierde ni se duplica por falta de señal. & Semanal, sin medición. & 0 pérdidas y 0 duplicados. & Prueba de desconexión de 14 horas y conciliación diaria en marcha blanca. \\
  R18-07 & Ruta del día siguiente en menos de 20 minutos, corregible y con registro. & 3,5 horas de una persona. & < 20 minutos. & Marcha blanca E1: medición diaria del tiempo de generación. \\
  R18-08 & POD digital disponible para el cliente el mismo día. & Guía en papel; 12 días para un reclamo. & 100 \% de entregas con POD el mismo día. & Marcha blanca E1: POD sobre entregas, más constancia de que el cliente la recibió o pudo consultarla el mismo día por un canal E1 (comprobante impreso o consulta asistida en la mesa de servicio), sin exigir dispositivo al canal tradicional. \\
  R18-09 & Cero guías extraviadas o ilegibles. & 1,1 \% mensual. & 0 \%. & Marcha blanca E1: verificar que todas las guías de despacho emitidas durante el período puedan localizarse y consultarse íntegramente en forma legible; cero guías extraviadas o ilegibles. \\
  R18-10 & Rendición cuadrada el mismo día con toda diferencia explicada. & Diferencias mensuales sin investigar ni asociar a una entrega. & 100 \% de diferencias con causal el mismo día. & Marcha blanca E1: cierre diario por camión. \\
  R18-11 & Costo de servir por cliente y entrega, desde el hecho. & Prorrateo por zona con criterio de 2016. & 100 \% de entregas con costo calculado. & Marcha blanca E2: reconciliación con la contabilidad del mes. \\
  R18-12 & Parque de envases controlado y pérdida anual no superior a la meta del SD2. & 14 \% estimado del parque al año, sin control. & Hasta 7 \% del parque al año (LafroX, 2026c, Tabla 2.1; S-26). & Aceptación provisional al cierre de la marcha blanca E1, con saldo conciliado por cliente y transportista; aceptación final tras 12 meses de operación. La pérdida anual se mide sobre el parque de envases, conforme al SD2, y no sobre la cantidad de despachos. \\
  R18-13 & Pedidos de cadenas por vía electrónica y aviso de despacho. & Cero pedidos electrónicos. & 100 \% de pedidos de cadenas certificadas por EDI. & Marcha blanca E2: pedidos EDI sobre pedidos de todas las cadenas comprometidas, certificadas antes del mes 21, y ASN emitidos dentro de 2 minutos sobre cargas (RF-12.09). La carga manual transitoria no acredita el resultado. \\
  R18-14 & La ocupación promedio aumenta y todos los camiones cumplen el mínimo comprometido al salir a reparto. & 68 \% promedio, con días de 41 \%. & Cada camión sale a reparto con una ocupación mínima del 60 \% de su capacidad volumétrica útil, incluidas las rutas rurales y sin excepciones; ocupación promedio superior al 68 \% (S-27). & Marcha blanca de la ola de reparto (4 semanas): verificar que cada salida cumpla el mínimo del 60 \%, sin excepciones, y que la ocupación promedio, calculada como volumen total cargado dividido por capacidad volumétrica útil total de las salidas evaluadas, sea superior al 68 \%. Comprobar además el cumplimiento de las restricciones de peso y condición térmica y de la promesa de entrega de S-21. \\
  R18-15 & La causa de cada faltante queda registrada al producirse. & Se registra el faltante, no la causa. & 100 \% de faltantes con causa registrada al producirse. & Marcha blanca E1: comprobar, mediante la trazabilidad del evento, que el 100 \% de los faltantes registre su causa al producirse, incluso sin conexión; la sincronización puede ser posterior. \\
  R18-16 & El planificador puede jubilarse sin que la operación se resienta. & El cumplimiento cae 6 a 9 puntos en sus vacaciones. & OTIF sin caída respecto del período con el planificador presente (propia). & Marcha blanca E1: dos semanas de planificación sin el planificador, comparadas con dos semanas previas de carga, rutas y mezcla de clientes comparables. Las diferencias atribuibles a demanda o incidencias externas se registran y excluyen, y las correcciones manuales quedan trazadas. \\
\end{tablalafrox}
